\section{Hierarchical-Memory and Virtual-Tensor Model}
\label{sec:model}

This section defines the objects manipulated by the compiler and runtime. The central abstraction is not a device tensor with occasional offload hooks, but a logical tensor whose physical materialization is a set of versioned regions distributed across a memory hierarchy.

\subsection{Tiers, paths, and capacities}

For each tier $\tier{k}$, let $C_k$ be its admitted byte capacity, $a_k$ its allocation granularity, and $q_k$ its alignment requirement. For every permitted directed transfer edge $e=(i,j)$, let $B_e(s)$ and $\lambda_e(s)$ denote calibrated effective bandwidth and latency for transfer size $s$, with concurrency limits $d_e$. A backend reports which transfers can overlap with compute and whether a direct path, such as storage-to-device DMA, exists. The planner never assumes that nominal bandwidth is achieved or that every pair of tiers is directly connected.

A concrete hierarchy might be
\[
\begin{aligned}
  \tier{0}&=\text{GPU HBM}, &
  \tier{1}&=\text{pinned host ring},\\
  \tier{2}&=\text{pageable DRAM cache}, &
  \tier{3}&=\text{NVMe files}.
\end{aligned}
\]
On a unified-memory system, logical tiers may share physical memory but retain different residency, pinning, and eviction policies. On CPU-only hardware, \tier{0} may be a bounded DRAM arena and \tier{1} a memory-mapped file.

\subsection{Virtual tensor}

\begin{definition}[Virtual tensor]
A virtual tensor $T$ is the descriptor
\[
T=(\mathbf n,d_{\ell},d_s,L,V,\chi,\rho),
\]
where $\mathbf n$ is logical shape, $d_{\ell}$ is logical dtype, $d_s$ is storage dtype or codec, $L$ is a logical layout and view mapping, $V$ is a versioned physical-region map, $\chi$ is integrity metadata, and $\rho$ is placement and recomputation policy.
\end{definition}

The logical index domain is
\[
  D_T=\prod_{r=0}^{k-1}\{0,\ldots,n_r-1\}.
\]
A \emph{tile} $\tau=(\mathbf o,\mathbf e)$ denotes a rectangular subset with origin $\mathbf o$ and extent $\mathbf e$. More general gather/scatter regions are represented by an index plan plus rectangular source spans. A materialization of $(T,\tau,v)$ is a physical buffer containing logical version $v$ of the tile in a declared layout and dtype.

The region map $V$ is many-to-many. One logical tile may have replicas in several tiers; one physical storage object may back multiple logical views or tied parameters. Every physical region is identified by a stable \code{storage\_id}, byte range, version, codec, checksum, and validity state. This prevents a model with tied input/output embeddings from being duplicated by the converter and preserves alias semantics through graph lowering.

\subsection{Tile states and coherence}

A tile instance transitions through the following abstract states:
\[
\begin{aligned}
\textsc{Absent}&\rightarrow\textsc{Reading}\rightarrow\textsc{HostReady}\\
&\rightarrow\textsc{Copying}\rightarrow\textsc{DeviceReady}\\
&\rightarrow\textsc{Computing}\rightarrow\textsc{Clean/Dirty}\\
&\rightarrow\textsc{Writing}\rightarrow\textsc{Absent}.
\end{aligned}
\]
Not every backend uses every state. Each transition is guarded by an event or future, and every buffer lease has an epoch. A buffer cannot be recycled until all events that reference its epoch have completed. This rule is what makes asynchronous deletion safe; Python reachability alone has no scheduling semantics.

Inference parameter tiles are immutable and may have multiple clean replicas. Mutable state uses single-writer or explicitly reduced multi-writer semantics. A dirty tile is durable only after its journal record and data checksum are committed according to the configured storage policy.

\subsection{Tensor classes}

The runtime distinguishes classes because their reuse and mutability differ:

\begin{enumerate}
  \item \emph{Parameters and constants} are immutable within an inference version and are candidates for prepacking and caching.
  \item \emph{Activation frontier} tensors are request-specific, usually short-lived, and may be spilled or rematerialized.
  \item \emph{Accumulators} hold partial reductions and have stricter numeric and writeback ordering.
  \item \emph{KV or recurrent state} grows with sequence length or evolves per token; it needs page allocation, sharing, copy-on-write, and request lifecycle management.
  \item \emph{Outputs and logits} may be streamed to downstream operators without full materialization.
  \item \emph{Gradients and optimizer state} are mutable, versioned, and transactionally updated in training.
  \item \emph{Transfer and kernel workspace} buffers are not model tensors but must be included in every memory proof.
\end{enumerate}

The invariant applies to all material classes. Pinning only $X$ and $Y$ while streaming weights is an optimization, not a valid universal assumption.

\subsection{Functional graph and stream plans}

Let $G=(N,E)$ be a finite functional graph after control-flow specialization for a request segment. Every node $n\in N$ is labeled by an operator $\omega\in\Oset$, logical input and output tensors, attributes, and semantic mode. Edges represent value dependencies, not physical buffers.

\begin{definition}[Stream-decomposable operator]
An operator $\omega$ is stream-decomposable under budget vector $\mathbf C$ if there exists a finite task graph $P_\omega$ such that:
\begin{enumerate}
  \item every task reads and writes finite regions of virtual tensors;
  \item each task declares a bounded resource vector and a primitive kernel supported by some backend;
  \item the peak admitted use of every tier does not exceed $\mathbf C$;
  \item the completed output has the declared semantics; and
  \item dependencies form a terminating schedule under the progress assumptions.
\end{enumerate}
The least capacity vector for which such a plan exists is the operator's \emph{minimum legal working set} for that backend and semantic mode.
\end{definition}

A stream plan is parameterized by shapes and may select among several loop nests. For example, a linear operator may keep an output accumulator stationary while sweeping input tiles, or keep an input tile resident while producing several output tiles. Both are correct; they differ in IO and memory.

\subsection{Resource vectors and arenas}

Every task $u$ has a resource vector
\[
  \mathbf r(u)=
  (r_{0}^{\mathrm{persistent}},r_{0}^{\mathrm{scratch}},
   r_{1}^{\mathrm{pinned}},r_{2}^{\mathrm{cache}},r_{3}^{\mathrm{spill}},
   s_{\mathrm{copy}},s_{\mathrm{compute}},f_{\mathrm{io}},\ldots).
\]
The memory broker owns preallocated or otherwise hard-bounded arenas for guaranteed resources. A task is dispatchable only when its full vector can be atomically leased. Kernel libraries with data-dependent or algorithm-selected workspaces must either accept a caller-provided workspace bounded by the plan, be profiled and locked to an algorithm, or be excluded from the guaranteed path.

For one device, a typical admitted bound is
\begin{equation}
\label{eq:device-bound}
M_0 \le R_0
 + q_W\lvert W_\tau\rvert
 + q_X\lvert X_\tau\rvert
 + q_Y\lvert Y_\tau\rvert
 + \lvert A_\tau\rvert
 + \lvert K_\tau\rvert
 + U_\tau + Z_0,
\end{equation}
where $R_0$ is runtime reserve, $q_*$ are in-flight buffer counts, $A_\tau$ is reduction state, $K_\tau$ is resident recurrent/KV state, $U_\tau$ is bounded kernel workspace, and $Z_0$ is alignment and metadata overhead. Total model parameter bytes do not occur in \eqref{eq:device-bound}.

\subsection{Plans as proof-carrying artifacts}

An execution plan is immutable and serializable. It contains:

\begin{itemize}
  \item graph and checkpoint content hashes;
  \item hardware and software capability fingerprints;
  \item shape guards and request envelope;
  \item operator implementation identifiers and semantic modes;
  \item tile grids, loop orders, buffer counts, and workspace maxima;
  \item per-tier memory and spill bounds;
  \item task dependency templates and priority classes;
  \item fallback chain and safe replan points;
  \item expected IO volume and calibrated cost estimate.
\end{itemize}

The runtime validates the plan against current reservations before use. The plan is ``proof-carrying'' in the engineering sense that every allocation and task maps to a statically checkable budget item; it is not a machine-checked proof unless an optional formal verifier is used.

\subsection{Architecture overview}
Figure~\ref{fig:architecture} summarizes the separation between functional semantics, planning, scheduling, and physical storage tiers.

\begin{figure}[H]
\centering
\resizebox{\linewidth}{!}{%
\begin{tikzpicture}[
  node distance=5mm and 8mm,
  box/.style={draw,rounded corners,align=center,minimum height=8mm,minimum width=24mm,fill=white},
  tierbox/.style={draw,rounded corners,align=center,minimum height=10mm,minimum width=31mm},
  arrow/.style={-{Latex[length=2.2mm]},thick}
]
  \node[box,fill=blue!8] (capture) {Graph capture\\and lowering};
  \node[box,fill=blue!8,right=of capture] (planner) {Tile planner\\and cost model};
  \node[box,fill=blue!8,right=of planner] (plan) {Validated\\execution plan};
  \node[box,fill=green!8,below=10mm of planner] (sched) {Credit scheduler, memory broker, leases, events};
  \node[tierbox,fill=orange!10,below left=10mm and -3mm of sched] (store) {Persistent store\\AMS / safetensors / mmap};
  \node[tierbox,fill=orange!10,right=of store] (host) {Bounded host cache\\and pinned ring};
  \node[tierbox,fill=orange!10,right=of host] (device) {Device arena\\and kernels};
  \node[box,fill=purple!8,below=8mm of host] (vt) {Virtual tensor map, versions, checksums, KV pages};
  \draw[arrow] (capture) -- (planner);
  \draw[arrow] (planner) -- (plan);
  \draw[arrow] (plan) -- (sched);
  \draw[arrow] (sched) -- (store);
  \draw[arrow] (sched) -- (host);
  \draw[arrow] (sched) -- (device);
  \draw[arrow,<->] (store) -- (host);
  \draw[arrow,<->] (host) -- (device);
  \draw[arrow,dashed] (vt) -- (store);
  \draw[arrow,dashed] (vt) -- (host);
  \draw[arrow,dashed] (vt) -- (device);
\end{tikzpicture}%
}
\caption{\AMS separates logical tensors and graph semantics from physical residency. The plan fixes bounded resources; the scheduler moves versioned tiles through the hierarchy.}
\label{fig:architecture}
\end{figure}
